\section{Supplementary Material}
\setcounter{figure}{0}
\renewcommand{\thefigure}{S\arabic{figure}}
\setcounter{table}{0}
\renewcommand{\thetable}{S\arabic{table}}

\subsection*{Forward-model characterization}

With increasing flexion under compression, the tibia of the nominal knee translated anteriorly relative to the femur, from \SIrange{2.0}{9.0}{\milli\meter} between \SI{0}{\degree} and \SI{90}{\degree}, and rotated internally from \SIrange{3}{16}{\degree} (\autoref{fig:forward}A).
The MCL carried most of the ligament force in extension, the PCL took over between \SI{65}{\degree} and \SI{70}{\degree} of flexion, and the LCL slackened within the first \SI{20}{\degree} (\autoref{fig:forward}B).
Across the \num{9999} converged characterization knees (the forward solution failed to converge for one of the \num{10000} knees), the median lateral gap exceeded the median medial gap at every flexion angle, and varus laxity exceeded valgus laxity.
This asymmetry is consistent with the three-bundle medial restraint and the more conforming medial dish.
The median total laxity was \SIrange{9}{10}{\milli\meter} for AP, \SIrange{4}{9}{\degree} for VV, and \SIrange{22}{34}{\degree} for IE.
External rotation exceeded internal rotation in flexion.
This asymmetry may reflect the simplified ligament geometry and posterolateral restraint of the model.

\begin{figure}[H]
\centering
\includegraphics{figS2_forward.pdf}
\caption{\textbf{The virtual knee exhibits flexion-dependent kinematics, ligament recruitment, and laxity.}
(A)~With increasing flexion under \SI{50}{\newton} compression, the tibia of the nominal knee translated anteriorly ($u_\mathrm{AP}$) and rotated internally ($\gamma$), with little varus rotation ($\beta$).
(B)~The MCL carried most of the neutral ligament force $F$ in extension, whereas the PCL carried most of it in deep flexion.
(C)~The lateral gap $g$ under \SI{100}{\newton} distraction per compartment exceeded the medial gap.
(D--F)~AP translation $\Delta u_\mathrm{AP}$ under \SI{100}{\newton}, VV rotation $\Delta\beta$ under \SI{10}{\newton\meter}, and IE rotation $\Delta\gamma$ under \SI{3}{\newton\meter} varied with flexion and between knees.
Positive values are anterior, varus, and internal.
In (C--F), lines show the median and bands the 5th to 95th percentiles of the \num{9999} converged characterization knees.
The nominal knee used the nominal attachment sites and the parameters of \autoref{tab:nominal}.}
\label{fig:forward}
\end{figure}

\begin{table}[H]
\centering
\caption{Ligament parameters of the nominal knee in \autoref{fig:forward}A,B and of the prior center. The prior center is the log-midpoint of the stiffness range and the midpoint of the reference-strain range in \autoref{tab:prior}. Stiffness is per bundle and per unit strain.}
\label{tab:nominal}
\begin{tabular}{llcccc}
\toprule
 & & MCL & LCL & PCL & Capsule \\
\midrule
Stiffness $k$ [\si{\kilo\newton}] & nominal knee & 2.75 & 2.00 & 4.50 & 1.50\\
 & prior center & 2.87 & 2.00 & 4.74 & 1.22\\
Reference strain $\varepsilon_\mathrm{r}$ [\si{\percent}] & nominal knee & 4 & 4 & 0 & 0\\
 & prior center & 4 & 4 & 0 & 0\\
\bottomrule
\end{tabular}
\end{table}

\subsection*{Equilibrium solver}

The Jacobian of \autoref{eq:equilibrium} was computed by forward differences, and each Newton step was limited to \SI{2}{\milli\meter} for AP and medial--lateral translation, \SI{1}{\milli\meter} for proximal translation, and \SIlist{0.05;0.1}{\radian} for varus and internal rotation.
Iterations stopped when the largest residual fell below \SI{1e-4}{\newton}, with moments divided by a length scale of \SI{25}{\milli\meter}.
If a solve did not reach a residual below \SI{0.01}{\newton}, the load was applied in 3, 8, and then 20 equal increments.

\subsection*{Path dependence of the equilibrium solutions}

In a check on 80 test knees, applying every load in 10 increments changed 2 of 2560 measurements by more than \SI{0.05}{\milli\meter} or \SI{0.05}{\degree}.
Both differences occurred under IE moments in flexion, where the model admitted two equilibria under the same load.
All other measurements changed by less than this threshold.
The same branches affected the Fisher analysis.
In 27 of the 499 knees, at least one central difference spanned two equilibria and produced an outlying sensitivity.
Excluding these knees changed the median identifiability indices by at most 0.02.

\subsection*{Robustness of the Fisher analysis}

A central-difference step of 0.001 instead of 0.02 on the first 50 knees changed the median identifiability indices by at most 0.02.
We also recomputed the indices from the same sensitivity matrices with all measurement-noise standard deviations doubled.
For the full examination, the mean of the eight median indices fell from 0.62 to 0.44.
However, the ranking of the nine measurement sets by this mean was unchanged, and for each parameter the Spearman rank correlation between the two noise levels was at least 0.98.

\subsection*{Variability of the held-out prediction errors}

\autoref{fig:calibration}B,C pools the held-out errors over knees.
\autoref{tab:heldout_iqr} gives the median and interquartile range of the per-knee errors.
The per-knee medians ranked the estimators as the pooled errors did, except that the gaps-only fits and the prior center had similar medians.

\begin{table}[H]
\centering
\caption{Median (interquartile range) across the 50 knees of the per-knee RMSE of the held-out predictions, each over nine outputs.}
\label{tab:heldout_iqr}
\begin{tabular}{lccc}
\toprule
Estimator & AP [\si{\milli\meter}] & VV [\si{\degree}] & IE [\si{\degree}] \\
\midrule
prior center & 1.30 (0.82--1.87) & 1.22 (0.93--1.64) & 3.11 (2.27--4.91)\\
gaps only, nominal attachment sites & 1.42 (0.90--1.81) & 1.13 (0.79--1.49) & 3.46 (2.22--4.52)\\
full examination, nominal attachment sites & 0.49 (0.39--0.71) & 0.58 (0.43--0.77) & 1.15 (0.86--1.59)\\
full examination, true attachment sites & 0.23 (0.18--0.29) & 0.26 (0.16--0.34) & 0.55 (0.40--0.68)\\
true parameters, nominal attachment sites & 0.97 (0.55--1.59) & 0.68 (0.43--1.05) & 2.28 (1.35--3.80)\\
\bottomrule
\end{tabular}
\end{table}

\subsection*{Stopping criteria of the least-squares fits}

All fits stopped before the budget of 500 residual evaluations. SciPy reports why each fit stopped: 1 for the gradient tolerance, 2 for the fit-error tolerance, 3 for the step tolerance, and 4 for both the fit-error and step tolerances. The 150 full-examination fits with nominal attachment sites ended with status 2, 3, and 4 in 7, 139, and 4 fits, after at most 68 evaluations. The fits with true attachment sites ended with status 2, 3, and 4 in 19, 118, and 13 fits, after at most 105 evaluations. The gaps-only fits ended with status 1, 2, and 3 in 14, 7, and 129 fits, after at most 222 evaluations. A stopping criterion establishes numerical termination, not global optimality.
